\section*{Extended Data Table 7 | Sensitivity of RNA test-set recovery to chemically similar reference actives}
Recovery is defined by Eqs.~(9)--(10) (ECFP4-Tanimoto window $0.25$--$0.50$). Analogue filters remove reference actives with max ECFP4-Tanimoto similarity $\geq0.5$ to the fine-tuning training partition (FT) or ligand-only pretraining library (PT). Each change compares filtered and original references on the identical eligible-target subset. The full model shows the largest reduction under the PT filter ($-0.198$; Pretrain+Constrain: $-0.099$; Pretrain Only: $-0.014$). Filtering removes reference actives only and does not affect the generated molecules.

\begin{center}\footnotesize
\setlength{\tabcolsep}{5pt}
\begin{tabular}{@{}lccc@{}}
\toprule
 & Unfiltered & \multicolumn{2}{c}{Change after analogue removal} \\
\cmidrule(lr){3-4}
Model configuration & Recovery (27 targets) & FT (18 targets) & PT (14 targets) \\
\midrule
No-pretraining baseline & 0.141 & $-0.077$ & --- \\
Pretrain Only & 0.034 & $-0.017$ & $-0.014$ \\
Pretrain+Constrain & 0.132 & $-0.016$ & $-0.099$ \\
\textbf{Full NA2smol} & \textbf{0.218} & $-0.077$ & $-0.198$ \\
\bottomrule
\end{tabular}
\end{center}
